\documentclass[11pt,a4paper]{article}
\usepackage{turkos}

\begin{document}
\phaseheader{12}{Storage: Block Devices, ATA Driver and Buffer Cache}{Make data survive a reboot: talk to a hard disk, organise devices behind one interface, and cache blocks in memory}{77}{83}

\ataglance{2 weeks}{4}{Phase 9: semaphores and mutexes; Phase 4: IRQ handling; Phase 7: the heap and lists}{An ATA PIO driver that identifies,
reads and writes the QEMU disk (first polling, then interrupt-driven), a block device layer with partitions from the
MBR, a RAM disk from a boot module, a buffer cache with LRU replacement and write-back, and \code{lsblk} and
\code{hexdump} commands. Data you write survives a reboot, and you can check it from Linux.}

\section{Why this phase}

Everything Turk-OS knows vanishes when QEMU exits. \OSTEP's third big idea, \textbf{persistence}, starts here: a
driver for a disk, and the layers that let the rest of the kernel use disks without knowing how they work. The file
system in Phase 13 sits directly on top of this phase and never touches hardware.

You will write the most complete driver so far. The ATA disk is controlled through a handful of I/O ports, it is slow
compared with the CPU, and it signals completion with an interrupt, so the task that asked for data should sleep
until it arrives. That brings together port I/O (Phase 2), interrupts (Phase 4) and blocking synchronisation
(Phase 9). Above the driver you will build the two layers that every UNIX kernel has in some form: a
device-independent \textbf{block layer} and a \textbf{buffer cache} that keeps recently used blocks in memory, because
reading from RAM is thousands of times faster than reading from a disk.

\section{Concepts you need}

\subsection{How a disk is organised}

A disk stores data in fixed-size \textbf{sectors} of 512 bytes. Old interfaces addressed them by cylinder, head and
sector (CHS); everything modern uses \textbf{logical block addressing} (LBA), which simply numbers the sectors 0, 1,
2, \dots\ \OSTEP chapter 37 builds a model of a mechanical disk's access time: \textbf{seek} time to move the arm,
\textbf{rotational} delay waiting for the sector to come round, and \textbf{transfer} time. The model explains why
sequential access is so much faster than random access, why disk schedulers reorder requests (FIFO, SSTF, SCAN,
C-SCAN), and why every file system tries to keep related blocks close together. QEMU's disk is a file on your host,
so none of the timing is real, but the interface is the same as on a real PC.

\subsection{The layers of I/O software}

\MOS\,\S5.3 divides I/O software into four layers, and Turk-OS follows it (Figure~\ref{fig:layers}). The
\textbf{interrupt handler} does the minimum and wakes whoever is waiting. The \textbf{device driver} knows one kind of
hardware. The \textbf{device-independent} layer (block devices, partitions, the buffer cache) gives every driver the
same interface and does the work common to all of them. \textbf{User-level} software (the file system, then
programs) only ever sees that interface.

\begin{figure}[htbp]
\centering
\begin{tikzpicture}[l/.style={draw=tkNavy!60, rounded corners=2pt, minimum width=8.8cm, text width=8.6cm, minimum height=0.62cm, font=\sffamily\scriptsize, align=center},
  n/.style={font=\sffamily\scriptsize, text=tkGray, anchor=west, align=left}]
  \node[l, fill=tkGreen!10] (u) at (0,0) {user process: \code{read(fd, buf, n)}};
  \node[l, fill=tkRed!8, below=2pt of u] (s) {system call layer (Phase 10)};
  \node[l, fill=tkBlue!8, below=2pt of s] (f) {VFS and TurkFS: file offset $\rightarrow$ block number (Phase 13)};
  \node[l, fill=tkAmber!14, below=2pt of f] (c) {\textbf{buffer cache}: \code{bread(dev, block)}, LRU, write-back};
  \node[l, fill=tkAmber!14, below=2pt of c] (b) {\textbf{block layer}: \code{struct blockdev}, partitions (hda1 = hda + offset)};
  \node[l, fill=tkAmber!14, below=2pt of b] (d) {\textbf{drivers}: ATA disk, RAM disk};
  \node[l, fill=tkAmber!14, below=2pt of d] (i) {\textbf{interrupt handler} (IRQ 14): acknowledge, wake the waiting task};
  \node[l, fill=tkGray!15, below=2pt of i] (h) {hardware: ATA controller and disk (\code{disk.img} on the host)};
  \draw[decorate, decoration={brace, amplitude=5pt}, tkAmber!80!black, line width=0.7pt] ([xshift=4pt]c.north east) -- ([xshift=4pt]b.south east)
    node[midway, right=6pt, n] {device-independent\\(this phase)};
  \draw[decorate, decoration={brace, amplitude=5pt}, tkAmber!80!black, line width=0.7pt] ([xshift=4pt]d.north east) -- ([xshift=4pt]i.south east)
    node[midway, right=6pt, n] {device-dependent\\(this phase)};
\end{tikzpicture}
\caption{The storage stack. A \code{read} on a file passes through every layer; a cache hit stops at the buffer cache.}
\label{fig:layers}
\end{figure}

\subsection{Programmed I/O on the ATA interface}

The primary ATA channel of a PC is a set of eight I/O ports starting at \code{0x1F0}, plus a control port at
\code{0x3F6} (Table~\ref{tab:ata}). To read, the driver selects the drive, writes the sector count and the 28-bit LBA
into four registers, and writes the \code{READ SECTORS} command. When the drive has a sector ready it sets
\textbf{DRQ} in the status register (and raises IRQ 14 if interrupts are enabled), and the driver reads 256 16-bit
words from the data port. That is \textbf{programmed I/O} (PIO): the CPU copies every word. \textbf{DMA}, where the
controller copies directly into memory, is faster and is a stretch goal (\OSTEP\,\S36.5, \MOS\,\S5.1.4).

\begin{table}[htbp]
\centering\small\renewcommand{\arraystretch}{1.2}
\begin{tabularx}{\linewidth}{@{}l >{\raggedright\arraybackslash}X >{\raggedright\arraybackslash}X@{}}
\toprule
\tkth{Port} & \tkth{Read} & \tkth{Write} \\
\midrule
\code{0x1F0} & data (16-bit) & data (16-bit) \\
\code{0x1F1} & error & features \\
\code{0x1F2} & sector count & sector count \\
\code{0x1F3}, \code{0x1F4}, \code{0x1F5} & LBA bits 0--7, 8--15, 16--23 & same \\
\code{0x1F6} & drive/head & \code{0xE0} $|$ (slave $\ll$ 4) $|$ LBA bits 24--27 \\
\code{0x1F7} & status: BSY \code{0x80}, DRDY \code{0x40}, DF \code{0x20}, DRQ \code{0x08}, ERR \code{0x01} & command: \code{0xEC} identify, \code{0x20} read, \code{0x30} write, \code{0xE7} flush \\
\code{0x3F6} & alternate status (reading it doesn't acknowledge an interrupt) & device control: bit 1 nIEN disables IRQs \\
\bottomrule
\end{tabularx}
\caption{The primary ATA channel's registers in 28-bit LBA mode.}
\label{tab:ata}
\end{table}

\subsection{The buffer cache}

The file system reads the same blocks over and over: the superblock, the bitmaps, the inodes of directories on every
path. The \textbf{buffer cache} keeps a fixed pool of block-sized buffers in memory. \code{bread(dev, block)} returns
the buffer for that block, reading it from disk only if it isn't cached; the caller holds the buffer's lock while
using it and gives it back with \code{brelse}. When all buffers are in use, the least recently used one that nobody
holds is reused. Writes can go to disk immediately (\textbf{write-through}) or be delayed until the buffer is evicted
or \code{sync} runs (\textbf{write-back}), which is faster but loses data in a crash: the central trade-off of Phase
13's crash-consistency discussion. The cache also guarantees there is only ever \emph{one} buffer per block, so two
tasks can't hold diverging copies.

\section{Reading plan}

\begin{readingplan}
\must & \OSTEP & Ch.~36 \emph{I/O Devices}, especially \S36.8 \emph{A Simple IDE Disk Driver} & 405--416 & Section 36.8 is essentially the driver of Tasks 12.2--12.4. \\
\must & \OSTEP & Ch.~37 \emph{Hard Disk Drives} & 419--436 & The disk model, access time, and disk scheduling. \\
\must & \MOS & \S5.2--5.3 \emph{Principles of I/O Software}, \emph{I/O Software Layers} & 351--369 & The layering this phase implements. \\
\must & \MOS & \S5.4 \emph{Disks} (hardware, formatting, arm scheduling, error handling, stable storage) & 369--388 & What a disk driver has to deal with in practice. \\
\should & \OSDI & \S3.4--3.7: I/O in MINIX 3, block devices, RAM disks, disks & 252--302 & A complete disk driver and RAM disk with explanations. \\
\should & \OSC & \S11.1--11.2 \emph{Mass-Storage Structure}, \emph{HDD Scheduling}; \S11.5 \emph{Storage Device Management} & 449--461, 463--467 & HDDs versus SSDs, scheduling, partitions and boot blocks. \\
\should & \OSC & \S12.3--12.5 \emph{Application I/O Interface}, \emph{Kernel I/O Subsystem}, \emph{Transforming I/O Requests} & 500--519 & Block versus character devices; buffering and caching in the kernel. \\
\could & \OSTEP & Ch.~38 \emph{RAID} & 437--456 & Combining disks for speed and reliability. \\
\could & \OSDI & \S5.6.5 \emph{The Block Cache} & 557--559 & The MINIX buffer cache, which you will meet again in Phase 13. \\
\end{readingplan}

The OSDev Wiki page \emph{ATA PIO Mode} is the practical reference for the register sequences, including the quirks
of real hardware that QEMU forgives.

\begin{minixbox}
The MINIX hard disk driver, \texttt{drivers/at\_wini/at\_wini.c (12100)}, is a user-space process that receives
read and write requests as messages. Its \code{w_identify} does what your Task~12.2 does, and \code{w_transfer} with
\code{com_out} issues commands and waits for the interrupt. All block drivers share the main loop in
\texttt{drivers/libdriver/driver.c (11000)}: that is the device-independent layer, done as a library. The RAM disk is
\texttt{drivers/memory/memory.c (11600)}. The block cache is part of the file-system server, in
\texttt{servers/fs/cache.c (22400)}, with \code{get_block} and \code{put_block} playing the roles of your
\code{bread} and \code{brelse}.
\end{minixbox}

\section{Build it}

\task{A disk for QEMU}
Create a raw disk image and attach it as the primary master; add a Makefile target so every run has it. Booting with
\code{-kernel} or from the CD image doesn't conflict, because QEMU attaches the CD to the secondary channel.

\begin{lst}{sh}{terminal}
qemu-img create -f raw build/disk.img 64M
qemu-system-i386 -kernel build/kernel.elf -serial stdio \
    -drive file=build/disk.img,format=raw,if=ide,index=0,media=disk
\end{lst}
\verify{The QEMU monitor's \code{info block} lists \code{ide0-hd0} backed by \code{build/disk.img}.}

\task{Find the drive with IDENTIFY}
Select the master drive (\code{0xA0} to \code{0x1F6}), zero the count and LBA registers, send \code{0xEC}, and read the
status. Zero means no drive. Otherwise wait for BSY to clear; if the LBA mid and high registers are now non-zero, the
device is not an ATA disk. Wait for DRQ (or ERR), then read 256 words. Words 27--46 hold the model name, with the two
bytes of every word swapped, and words 60--61 hold the number of 28-bit addressable sectors.
\verify{At boot the kernel prints \code{hda: QEMU HARDDISK, 131072 sectors (64 MiB)}.}

\task{Read and write sectors, polling}
Write \code{ata_read} and \code{ata_write} for up to 255 sectors per call. After selecting the drive, wait about
400\,ns by reading the alternate status port four times, because the status is not valid immediately. Before every
sector, wait until BSY is clear and DRQ is set, and check ERR and DF. Move data with \code{rep insw} and
\code{rep outsw}. After writing, send \code{CACHE FLUSH} (\code{0xE7}) and wait for BSY to clear, or the data may still
be sitting in the drive's cache.

\begin{lst}{c}{kernel/drivers/ata.c}
#define ATA_DATA 0x1F0
#define ATA_SECCOUNT 0x1F2
#define ATA_LBA0 0x1F3
#define ATA_LBA1 0x1F4
#define ATA_LBA2 0x1F5
#define ATA_DRIVE 0x1F6
#define ATA_CMD 0x1F7
#define ATA_STATUS 0x1F7
#define ATA_ALTSTATUS 0x3F6
#define ST_BSY 0x80
#define ST_DF 0x20
#define ST_DRQ 0x08
#define ST_ERR 0x01

static inline void insw(uint16_t port, void *buf, uint32_t words)
{
    __asm__ volatile ("rep insw" : "+D"(buf), "+c"(words) : "d"(port) : "memory");
}

static void delay400ns(void) { for (int i = 0; i < 4; i++) inb(ATA_ALTSTATUS); }

static int wait_ready(void)
{
    uint8_t s;
    while ((s = inb(ATA_STATUS)) & ST_BSY)
        ;
    if (s & (ST_ERR | ST_DF))
        return -EIO;
    return (s & ST_DRQ) ? 0 : -EIO;
}

int ata_read(uint32_t lba, uint8_t count, void *buf)
{
    outb(ATA_DRIVE, 0xE0 | ((lba >> 24) & 0x0F));   /* LBA mode, master, LBA 24..27 */
    delay400ns();
    outb(ATA_SECCOUNT, count);
    outb(ATA_LBA0, lba & 0xFF);
    outb(ATA_LBA1, (lba >> 8) & 0xFF);
    outb(ATA_LBA2, (lba >> 16) & 0xFF);
    outb(ATA_CMD, 0x20);                            /* READ SECTORS */
    for (uint8_t i = 0; i < count; i++) {
        if (wait_ready() < 0)
            return -EIO;
        insw(ATA_DATA, (uint8_t *)buf + i * 512, 256);
    }
    return 0;
}
\end{lst}
\verify{Write a pattern to sectors 100--107, read it back and compare. Then exit QEMU and look at the image on
the host with \code{hexdump -C -s $((100*512)) -n 64 build/disk.img}: your pattern is there.}

\task{Let the disk interrupt}
Polling wastes the CPU while the disk works. Enable disk interrupts (clear nIEN in the device control register),
unmask IRQ 14 on the slave PIC and IRQ 2 on the master, and register a handler for vector 46 that reads the status
register (which acknowledges the interrupt) and calls \code{sem_up}. In \code{ata_read}, replace the busy wait before
each sector with \code{sem_down}. Serialise requests with a mutex, since the controller can only do one thing at a
time.
\verify{While one task reads 10\,MiB from the disk, another task that prints once a second keeps its rhythm, and
\code{ps} shows the reader blocked most of the time.}

\task{A block device layer}
Define \code{struct blockdev} with a name, the sector size, the number of sectors, \code{read} and \code{write}
function pointers, and a private pointer for the driver. Keep a registry, and add \code{blockdev_find("hda")}. From
now on, nothing above this layer calls \code{ata_read} directly.

\begin{lst}{c}{include/blockdev.h}
struct blockdev {
    char     name[8];                   /* "hda", "hda1", "ram0" */
    uint32_t sector_size;               /* 512                   */
    uint32_t nsectors;
    int    (*read)(struct blockdev *d, uint32_t lba, uint32_t count, void *buf);
    int    (*write)(struct blockdev *d, uint32_t lba, uint32_t count, const void *buf);
    void    *priv;                      /* driver state; for a partition: parent + offset */
};
\end{lst}

\task{Partitions}
Read sector 0 of each disk. If bytes 510--511 are \code{55 AA}, parse the four 16-byte partition entries at offset
\code{0x1BE}: byte 4 is the type, bytes 8--11 the first LBA, bytes 12--15 the number of sectors. Register each used
entry as a block device (\code{hda1}\dots\code{hda4}) whose read and write add the offset and reject any request that
would cross the partition's end. Partition the image on the host; type \code{81} is the traditional MINIX type, which
fits the file system of Phase 13.

\begin{lst}{sh}{terminal}
echo 'start=2048, type=81' | sfdisk build/disk.img
\end{lst}
\verify{\code{lsblk} lists \code{hda} (64\,MiB) and \code{hda1} starting at sector 2048, and a read beyond the end of
\code{hda1} fails with \code{-EINVAL}.}

\task{A RAM disk}
Register any boot module named \code{ramdisk} as the block device \code{ram0}; its read and write functions are just
\code{memcpy}. It is useful for testing the file system without the ATA driver in the way, and in Phase 13 the initial
root file system can live there.

\task{The buffer cache}
Write \code{bread}, \code{bwrite}, \code{bdirty}, \code{brelse} and \code{bsync} over a pool of 64 buffers of
1\,KiB (the block size of the Phase 13 file system, two sectors). Each buffer has a device, a block number, valid and
dirty flags, a reference count, a mutex held by whoever uses the data, and a position in an LRU list; add a small hash
table from (device, block) to buffer once the linear search shows up in profiles. Look up and claim buffers under one
spinlock, so two tasks asking for the same block get the same buffer.

\begin{lst}{c}{include/bcache.h}
#define BLOCK_SIZE 1024

struct buf {
    struct blockdev *dev;
    uint32_t         blockno;
    bool             valid;             /* data holds the block's contents    */
    bool             dirty;             /* data must be written back          */
    int              refcnt;            /* users holding it (not reusable)    */
    struct mutex     lock;              /* held while a task uses data        */
    struct list_node lru;               /* most recently released at the head */
    uint8_t          data[BLOCK_SIZE];
};

struct buf *bread(struct blockdev *dev, uint32_t blockno);   /* returns locked, valid */
void        bdirty(struct buf *b);                            /* write back later       */
void        bwrite(struct buf *b);                            /* write now              */
void        brelse(struct buf *b);                            /* unlock, move to MRU    */
void        bsync(void);                                      /* flush every dirty one  */
\end{lst}
\verify{Reading the same block twice causes one disk read; a \code{cache} command prints hits, misses and dirty
buffers; after \code{bsync} the dirty count is zero.}

\task{Tools and tests}
Add \code{lsblk} and \code{hexdump <dev> <block>} to the monitor. Add tests for block device reads and writes on
\code{ram0}, partition bounds checking, cache hits and eviction (touch 100 different blocks with 64 buffers), and a
persistence test: write a known block, \code{bsync}, reboot with the same image, and check it.

\section{Testing and debugging}

\begin{pitfalls}
IDENTIFY returns only zeros or hangs & Wrong drive select, or no disk attached at \code{index=0} & Check \code{info block}; select with \code{0xA0} and wait 400\,ns before reading status. \\
Reads return the previous sector's data & DRQ not checked before each sector of a multi-sector read & Call \code{wait_ready} before every 256-word transfer. \\
Written data is missing after QEMU exits & No \code{CACHE FLUSH} after writes, or dirty buffers never synced & Send \code{0xE7} after writing; call \code{bsync} before shutdown. \\
IRQ 14 never arrives & IRQ 2 on the master or IRQ 14 on the slave still masked, or nIEN set & Unmask both and clear nIEN in \code{0x3F6}; \code{info pic} shows the masks. \\
Two tasks see different contents for the same block & Two buffers were created for the same block & Do lookup and allocation atomically under the cache's spinlock. \\
Partition data lands in the wrong place & The partition offset is applied twice, or not at all & Test with a known pattern at the partition's first sector and check with \code{hexdump -C} on the host. \\
\end{pitfalls}

\section{Milestone}

\begin{milestonebox}[Data survives]
\begin{checklist}
  \item The ATA driver identifies the QEMU disk, and reads and writes with interrupts and a blocking wait.
  \item All disk access goes through \code{struct blockdev}; partitions from the MBR appear as their own devices.
  \item The RAM disk works from a boot module.
  \item The buffer cache serves repeated reads from memory and writes dirty blocks back on \code{bsync}.
  \item Data written by Turk-OS survives a reboot and is visible with \code{hexdump -C} on the host.
\end{checklist}
\end{milestonebox}

\begin{questionbox}
\begin{enumerate}
  \item What are the three components of disk access time, and which one does sequential access avoid?
  \item Compare programmed I/O and DMA. What does the CPU do during each?
  \item Why must the driver check DRQ before every sector rather than once per command?
  \item What does the device-independent block layer buy you? Give two concrete examples in Turk-OS.
  \item What does the buffer cache guarantee besides speed? What is the risk of write-back caching?
  \item Describe the layout of an MBR and one partition entry.
\end{enumerate}
\end{questionbox}

\section{Going further}

Enumerate the PCI bus through configuration ports \code{0xCF8}/\code{0xCFC}, find the IDE controller, and use its
bus-master DMA engine so the CPU no longer copies every word. Queue requests and sort them with C-SCAN
(\OSTEP\,\S37.5, \MOS\,\S5.4.3). Add read-ahead: when block $n$ is read, start reading $n+1$ in the background. Or write
a \code{virtio-blk} driver, which is how modern virtual machines expose disks, and compare its speed with ATA.

\nextphase{13}{File Systems: VFS, initrd and TurkFS}
\booklegend
\end{document}
